\section{La metodología de investigación de He Kaiming (何凯明)}

Una de las partes más valiosas de la entrevista es que Xie Saining (谢赛宁) dedica una gran cantidad de tiempo a desglosar el método de investigación de He Kaiming. Estos contenidos son importantes no porque ofrezcan una especie de «receta del éxito», sino porque reducen la investigación de punta, de un talento misterioso, a una forma de trabajo que puede observarse, entrenarse y también malinterpretarse. El juicio más central es: \textit{``Research nunca es un desarrollo lineal; un research que se desarrolla de forma lineal nunca es un buen research''}. Esta frase casi puede tomarse como el compendio de cómo él entiende el trabajo creativo. Una investigación verdaderamente buena no avanza en línea recta de la idea A al paper B; se desvía continuamente en el camino, se topa con obstáculos, se ve obligada a reescribir el problema, y el resultado final suele diferir mucho de lo que se planteó al inicio.

Esto también explica por qué él añade de inmediato: \textit{``la idea que se te ocurre al principio no es tu idea; la idea que surge durante la exploración sí es tuya''}. La ocurrencia que surge nada más sentarse suele ser solo un préstamo, ensamblado a partir de lecturas recientes, tendencias de moda o prejuicios ya existentes; solo después de empezar de verdad a experimentar, de encontrarse con fracasos y de ver señales extrañas, el problema empieza poco a poco a tomar un contorno propio. Xie Saining insiste mucho en el valor de la exploración, no porque la exploración tenga algo de romántico, sino porque solo ella puede forzar a salir a la luz la relación real entre uno mismo y el problema. Cuando alguien insiste siempre, sobre el papel, en que «ya tengo la respuesta completa», con frecuencia eso indica más bien que todavía no ha entrado de verdad en su objeto de estudio.

\begin{importantbox}{La no linealidad no es un efecto secundario, sino la naturaleza misma de la investigación}
Xie Saining no dice que la investigación a veces tome desvíos, sino que el «desvío» en sí mismo es la parte más informativa de la investigación. La experiencia de pasar meses sin ningún rumbo y de que todo se destrabe de pronto en el último mes es, a su juicio, precisamente la forma normal. Porque los hallazgos verdaderamente importantes suelen venir de que uno se ve obligado a cambiar de perspectiva frente a un obstáculo, y no de avanzar en línea recta siguiendo el plan original.
\end{importantbox}

Sobre esta comprensión no lineal, He Kaiming ha desarrollado además un ritmo de investigación bastante claro. Xie Saining lo resume como un ciclo de seis meses: durante el primer periodo, de uno a dos meses, se dedica sobre todo a la exploración abierta, para encontrar cualquier signal que valga la pena seguir; en los siguientes dos a tres meses se hace scale up de la dirección más prometedora, se sistematiza, se hacen comparaciones y depuración de errores; y en el último mes se entra en la escritura y el pulido. Este ritmo es eficiente no por la distribución del tiempo en sí, sino porque parte de la base de que «el problema cambiará durante la exploración». Por eso, en la etapa inicial el plan no puede fijarse de manera demasiado rígida, y en la etapa final tampoco se puede forzar una historia cuando no hay señal.

\begin{knowledgebox}{Un ciclo típico de investigación de seis meses}
\begin{tabularx}{\textwidth}{>{\raggedright\arraybackslash}p{0.18\textwidth} Y}
\toprule
Etapa & Foco de trabajo \\
\midrule
Meses 1--2 & Exploración intensiva, construcción de baseline, prueba y error rápidos, búsqueda de puntos anómalos y de signal útil, sin apresurarse a definir el método final. \\
Meses 3--5 & Organización de recursos en torno al signal más fuerte, scale up, experimentos de ablación, análisis de errores e ingeniería de sistemas necesaria, para comprimir los fenómenos ocasionales en resultados estables. \\
Mes 6 & Terminar por adelantado la redacción del cuerpo del paper, y luego someter repetidamente figuras, tablas, redacción y lógica narrativa a un proceso de polish, para que el paper presente una línea clara de avance del problema. \\
\bottomrule
\end{tabularx}
\end{knowledgebox}

Él subraya en particular un juicio contraintuitivo: el peor research no suele ser el que fracasa, sino aquel que de principio a fin no encuentra ningún obstáculo. Porque si la idea inicial coincide por completo con el resultado final, eso significa que probablemente uno nunca se acercó a lo verdaderamente desconocido, sino que solo ejecutó un plan demasiado conservador y de muy baja información. En cambio, un resultado negativo puede ser, más bien, una señal muy buena. Xie Saining pone como ejemplo que si cierta dirección hace «perder diez puntos», eso al menos indica que el sistema es muy sensible a ese factor, y que en sentido contrario tal vez haya una ganancia enorme; lo verdaderamente preocupante es que el rendimiento se quede estancado, ni bien ni mal, lo cual indica que uno todavía no ha tocado el grado de libertad clave del sistema.

\begin{warningbox}{El peor research: sin obstáculos}
Lo más impactante de la entrevista es que él define «no tener obstáculos» como la marca de una mala investigación. Porque no tener obstáculos suele significar que no se ha entrado de verdad en el núcleo del problema, que no se ha topado con una tensión estructural, y que tampoco se ha forzado la aparición de una explicación nueva. Por el contrario, la confusión, los resultados anómalos y los experimentos fallidos suelen estar más cerca de un hallazgo importante que un avance sin tropiezos.
\end{warningbox}

Esta metodología no tiene solo un «sabor filosófico»; también se expresa en una disciplina experimental muy estricta. He Kaiming exige que, antes de ejecutar cada experimento, el investigador escriba primero su propia predicción sobre el resultado. En apariencia, este gesto solo busca evitar el diagnóstico retrospectivo fácil, pero en realidad entrena al investigador para construir un modelo causal: ¿por qué crees que este cambio va a subir, a bajar o a quedarse igual? Si el resultado no coincide con la predicción, ¿el problema está en que la intuición se equivocó, o en que hay una variable en el sistema que todavía no se ha identificado? Xie Saining también menciona el uso de tablas de Excel para dar seguimiento a los experimentos; esta forma de gestión, en apariencia ordinaria, en el fondo combate la desviación cognitiva más común en el proceso de investigación. El cerebro humano tiende con facilidad a recordar solo los resultados que respaldan su propia narrativa, y a pasar por alto la distribución real de todo el proceso de búsqueda.

Un nivel más profundo es el llamado research taste. He Kaiming le regaló el Sutra del Diamante como obsequio de bienvenida al puesto, y esa frase, \textit{``Toda imagen es ilusoria; si ves que las imágenes no son la esencia, entonces ves al Tathagata''}, es tomada directamente por Xie Saining para explicar la estética de la investigación. El llamado taste no consiste en perseguir los términos más rimbombantes, sino en poder atravesar la «imagen» para ver la «esencia», y cuestionar qué variable esencial hay en verdad detrás de un paper, un efecto o una dirección de moda. Esto también explica por qué después siempre logra comprimir trabajos aparentemente complejos en unos cuantos juicios afilados, por ejemplo que ResNeXt es en esencia como un MoE, o que soft attention tal vez no sea la parte más importante de ViT. El gusto viene de identificar la esencia, no de acumular terminología.

\begin{importantbox}{El paper es como una película, la investigación es como un descubrimiento narrativo}
Xie Saining, apoyándose en el interés compartido con He Kaiming por la escritura cinematográfica, propone una comparación muy gráfica: la peor película es un recuento cronológico sin más, y el peor paper también lo es. El punto central de un buen paper no está solo en la technique, sino en «cómo llegaste hasta aquí en realidad». El problema, el conflicto, el giro, la evidencia y la solución final forman el arco narrativo de un paper. Aquí, el storytelling no es un empaque, sino organizar la estructura real del descubrimiento en una forma que el lector pueda entender.
\end{importantbox}

Esto también explica por qué He Kaiming puede terminar de escribir el paper un mes antes del deadline, y después pulirlo palabra por palabra y signo de puntuación por signo de puntuación, incluso con un OCD tal que ninguna línea de texto puede ocupar menos del 60\%. Este polish extremo no es formalismo, sino responsabilidad hacia la narrativa de la investigación. Un paper ambiguo, inflado o que se repite a sí mismo suele revelar que el propio autor todavía no tiene clara la verdadera estructura ósea del problema. Que Xie Saining ponga en paralelo «Story», del teórico de cine McKee (麦基), con la escritura de investigación, demuestra precisamente su convicción: tanto en la ciencia como en el arte, lo que de verdad conmueve no es la acumulación de material, sino si la estructura tiene fuerza.

\subsection{Resumen de la sección}

La metodología de He Kaiming puede resumirse en cuatro capas: reconocer la naturaleza no lineal del research, generar una idea verdadera mediante la exploración y no mediante la mera especulación, usar la disciplina experimental para aproximarse a una explicación causal, y luego usar la capacidad narrativa para expresar la estructura del hallazgo. Para Xie Saining, este método no es un principio abstracto, sino la gramática de trabajo detrás de todas sus obras representativas posteriores.
